\chapter{Broken Symmetry and Further Directions}

\section{Spontaneous symmetry breaking}

Consider a real scalar field with potential
\[
  V(\phi)=\frac{\lambda}{4}(\phi^2-v^2)^2,
  \qquad \lambda>0.
\]
The minima satisfy
\[
  \frac{\dd V}{\dd\phi}
  =
  \lambda(\phi^2-v^2)\phi=0.
\]
The points \(\phi=\pm v\) are minima, while \(\phi=0\) is a maximum if \(v^2>0\).  The Lagrangian is symmetric under \(\phi\to-\phi\), but choosing one vacuum, say \(\phi=v\), hides the symmetry.

Write
\[
  \phi(x)=v+h(x).
\]
Then
\[
  V=\frac{\lambda}{4}\left((v+h)^2-v^2\right)^2
  =
  \frac{\lambda}{4}(2vh+h^2)^2
  =
  \lambda v^2h^2+\lambda vh^3+\frac{\lambda}{4}h^4.
\]
Since a mass term appears as \(\frac12m_h^2h^2\) in \(V\),
\[
  m_h^2=2\lambda v^2.
\]
The fluctuation \(h\) is massive.

\section{Goldstone's theorem}

For a complex scalar with
\[
  V(\Phi)=\lambda\left(\Phi^*\Phi-\frac{v^2}{2}\right)^2,
\]
choose
\[
  \Phi(x)=\frac1{\sqrt2}(v+h(x)+\ii\chi(x)).
\]
Then
\[
  \Phi^*\Phi-\frac{v^2}{2}
  =
  vh+\frac12(h^2+\chi^2).
\]
The quadratic part of the potential is
\[
  V_{\text{quad}}=\lambda v^2h^2.
\]
There is no \(\chi^2\) term.  The field \(\chi\) is massless.  Goldstone's theorem says that every spontaneously broken continuous global symmetry gives a massless mode, under broad assumptions.

\section{The Higgs mechanism}

Now make the \(U(1)\) symmetry local:
\[
  \Lag=
  -\frac14F_{\mu\nu}F^{\mu\nu}
  +(D_\mu\Phi)^*(D^\mu\Phi)
  -\lambda\left(\Phi^*\Phi-\frac{v^2}{2}\right)^2.
\]
Use a gauge choice in which
\[
  \Phi(x)=\frac1{\sqrt2}(v+h(x)).
\]
Then
\[
  D_\mu\Phi=(\p_\mu+\ii qA_\mu)\frac{v+h}{\sqrt2}.
\]
The part quadratic in fields includes
\[
  \frac12q^2v^2A_\mu A^\mu.
\]
This is a mass term for the gauge field:
\[
  m_A=qv.
\]
The would-be Goldstone field has become the longitudinal polarization of a massive vector boson.  Gauge symmetry was not truly broken as a redundancy; rather, the spectrum reorganized around a vacuum with nonzero field value.

\section{Non-Abelian symmetry breaking}

For a scalar multiplet \(\Phi\) in a representation of a group \(G\), the covariant derivative is
\[
  D_\mu\Phi=\p_\mu\Phi+\ii gA_\mu^aT^a\Phi.
\]
If the vacuum expectation value is \(\Phi_0\), then the gauge-boson mass terms come from
\[
  (D_\mu\Phi_0)^\dagger(D^\mu\Phi_0)
  =
  g^2 A_\mu^aA^{b\mu}
  (T^a\Phi_0)^\dagger(T^b\Phi_0).
\]
Generators satisfying
\[
  T^a\Phi_0=0
\]
remain unbroken and their gauge bosons stay massless.  Generators that move the vacuum produce massive gauge bosons.  This is the algebraic heart of electroweak symmetry breaking.

\section{Anomalies}

A symmetry of the classical action may fail after quantization.  Such a failure is an anomaly.  For a massless Dirac fermion, the axial current
\[
  j_5^\mu=\bar\psi\gamma^\mu\gamma^5\psi
\]
is classically conserved.  In QED, quantum effects give
\[
  \p_\mu j_5^\mu
  =
  \frac{q^2}{16\pi^2}
  \epsilon^{\mu\nu\rho\sigma}F_{\mu\nu}F_{\rho\sigma}
  \]
for one Dirac fermion, with convention-dependent signs.  A derivation can be made by regulating the fermion measure in the path integral or by evaluating the triangle diagram.  The important lesson is that regularization is not neutral: preserving one symmetry may expose the failure of another.

Gauge anomalies are dangerous because gauge symmetry removes unphysical degrees of freedom.  A consistent gauge theory must have its gauge anomalies cancel.  This requirement strongly constrains the matter content of the Standard Model.

\section{Finite temperature}

Thermal quantum field theory replaces vacuum expectation values by thermal traces:
\[
  \avg{\mathcal O}_\beta
  =
  \frac{\Tr(e^{-\beta H}\mathcal O)}{\Tr(e^{-\beta H})},
  \qquad
  \beta=\frac1T.
\]
In the Euclidean path integral, time becomes imaginary and periodic:
\[
  \tau\sim\tau+\beta.
\]
Bosonic fields are periodic,
\[
  \phi(\tau+\beta,\bm x)=\phi(\tau,\bm x),
\]
while fermionic fields are antiperiodic,
\[
  \psi(\tau+\beta,\bm x)=-\psi(\tau,\bm x).
\]
The continuous energy integral becomes a sum over Matsubara frequencies.  This formalism describes phase transitions, plasmas, and early-universe field theory.

\section{Curved spacetime}

In curved spacetime the metric becomes position dependent:
\[
  \eta_{\mu\nu}\to g_{\mu\nu}(x).
\]
The invariant volume element is
\[
  \dd^4x\sqrt{-g},
\]
where \(g=\det g_{\mu\nu}\).  A scalar field action becomes
\[
  S=\int\dd^4x\sqrt{-g}
  \left[
  \frac12g^{\mu\nu}\p_\mu\phi\p_\nu\phi
  -\frac12m^2\phi^2
  \right].
\]
Quantization in curved backgrounds raises new questions because the notion of particle depends on the observer and on the choice of time coordinate.  Even so, the local field-theory tools developed in flat spacetime remain the starting point.

\section*{Exercises}
\addcontentsline{toc}{section}{Exercises}

\begin{exercise}
For \(V(\phi)=\lambda(\phi^2-v^2)^2/4\), compute \(V''(v)\).
\begin{answer}
\(V'=\lambda(\phi^2-v^2)\phi\), so \(V''=\lambda(3\phi^2-v^2)\), and \(V''(v)=2\lambda v^2\).
\end{answer}
\end{exercise}

\begin{exercise}
In the Abelian Higgs model, identify the term that gives the gauge boson a mass.
\begin{answer}
From \(|D_\mu\Phi|^2\) with \(\Phi=(v+h)/\sqrt2\), the quadratic term is \(\frac12q^2v^2A_\mu A^\mu\), so \(m_A=qv\).
\end{answer}
\end{exercise}

